\chapter{When it breaks}
\label{ch:breaks}

\section*{What this chapter is for}

Something will not work, in front of somebody, with a clock running. What is
assessed at that moment is not whether it broke but how long you spend on it
and whether you decide or flounder. A decision tree written in advance turns
that into an execution rather than an improvisation.

\section{The budget rule}

Every failure gets a time budget, decided before you start, and when it runs
out you take the fallback whether or not you feel close.

\judgement{``I am nearly there'' is the most expensive sentence in a timed
build, because it is true roughly half the time and the other half it is the
beginning of twenty minutes. You cannot tell which from the inside, which is
exactly why the budget is set in advance rather than judged at the time.}

\section{The tree}

\begin{kittable}{@{}Xl>{\raggedright\arraybackslash}p{22mm}X@{}}
\textbf{What broke} & \textbf{Budget} & \textbf{Fallback} & \textbf{Say} \\
\midrule
Tooling will not install & 2 min & do without it & ``not worth the time, going without'' \\
Database will not start & 3 min & in-memory store & ``switching to in-memory; it stops being evidence about concurrency, and I will say so'' \\
Container will not build & 3 min & run the process directly & ``running it directly; containerising is not what is being assessed'' \\
External dependency is down & 2 min & a stub, named as a stub & ``stubbing this; the mapping is untested and that is a known gap'' \\
Credential is missing & 2 min & the path that needs it is out of scope & ``I cannot exercise this path; here is what it would have proved'' \\
A library does not behave as expected & 5 min & do it by hand & ``writing this by hand rather than debugging the abstraction'' \\
Assistant is producing nonsense & 2 prompts & write it yourself & ``I will write this one'' \\
The design is wrong & --- & change it, out loud & ``this boundary does not work, here is why, here is what I am doing instead'' \\
\end{kittable}

The last row has no budget because it is not a failure. Discovering that the
design is wrong is the build doing its job, and Chapter~\ref{ch:design} put the
sentence in place for it in advance.

\section{The sentence matters as much as the fallback}

Every row has a \emph{say} column, and it is not decoration. A silent fallback
looks like the thing that was supposed to happen. A narrated one is a decision,
with a stated cost, which is what the round is looking for.

The structure is always the same three parts: what is not working, what you are
doing instead, and \textbf{what that costs}. The third is the one people leave
out, and it is the one that separates a fallback from a shortcut.

\begin{kitnote}
``Switching to in-memory'' is a shortcut. ``Switching to in-memory; note that
the concurrency test now proves nothing, because this provider has no
transactions'' is a decision with its consequence named. The second is a much
stronger thing to say, and it takes four extra words.
\end{kitnote}

\section{What not to do}

\textbf{Do not debug in silence.} Thirty seconds is fine. Two minutes of
silent scrolling with somebody watching is dead time for both of you, and the
fix is narration: what you think is wrong, what you are checking, what would
confirm it.

\textbf{Do not restart from scratch.} Under a clock it feels clean and it is
never right. The second attempt is slower, because the first one had
accumulated context.

\textbf{Do not disable the check.} A failing assertion removed to make a build
green is the one action in this round that changes an interviewer's assessment
of your judgement rather than your luck. If a check is wrong, say it is wrong
and say why.

\textbf{Do not apologise repeatedly.} Once is human. Beyond that it consumes
the time you were apologising for.

\section{Recovering the session}

If a failure has cost ten minutes and the run sheet is compromised, do not
carry on as though it has not.

\begin{kitmethod}
\textbf{1. Say where you are against the plan.} ``I have lost about ten minutes
on the database; I am cutting the second entity.'' Out loud, once.

\textbf{2. Re-cut to the next checkpoint only.} Do not re-plan the whole
session mid-session. One checkpoint ahead.

\textbf{3. Protect the demonstration block.} It is the last thing to be given
up, not the first. A session that ends mid-edit produces nothing to assess.
\end{kitmethod}

\judgement{A candidate who loses fifteen minutes to a genuine infrastructure
failure, says so, cuts scope deliberately, and demonstrates something smaller
on time has shown more of what the round is looking for than one whose session
ran smoothly \dash{} because the second one was never observed making a
decision under pressure.}

\begin{drill}
During the rehearsal in Chapter~\ref{ch:rehearsal}, have your partner break
something without warning: disconnect the network at minute forty, or ask you
to switch databases. The instruction for you is the hard one \dash{} do not
stop the rehearsal. The whole purpose is to find out what you do, and you
cannot find that out in a session that is halted when it gets uncomfortable.
\end{drill}

\section*{What this chapter does not claim}

The budgets are the author's and are deliberately short. The argument for short
budgets is that the cost of an unnecessary fallback is small and visible, and
the cost of an overrun is large and invisible until it is too late. That is a
judgement about asymmetry, not a measurement.

No claim is made about how failures are actually weighted by interviewers. What
is reported elsewhere is that reasoning is watched; this chapter assumes that
extends to reasoning under failure.
